% --------------------------------------------------
\subsubsection{Нэг хүчин зүйлийн зорилтот матриц}
% --------------------------------------------------

Нэг хүчин зүйлийн зорилтот матриц нь хөрөнгүүдийн өгөөжийн
хоорондын хамаарлыг зах зээлийн нийтлэг хүчин зүйлээр
тайлбарлах нэг хүчин зүйлийн загварт (Single-Index Model)
суурилдаг. Уг загварт хөрөнгийн өгөөжийг зах зээлийн
нийтлэг хүчин зүйлтэй холбоотой хэсэг болон тухайн
хөрөнгийн өөрийн онцлогоос хамаарах үлдэгдэл хэсэгт
хуваан авч үздэг \citep{ledoitwolf2003}.

Энэхүү судалгаанд зах зээлийн нийтлэг хүчин зүйлийг
S\&P 500 индексийн өгөөжөөр төлөөлүүлнэ.

$i$-р хөрөнгийн $t$-р хугацаан дахь өгөөжийг
\[
y_{ti}
=
a_i+\beta_i R_{m,t}+\varepsilon_{ti}
\]
гэж загварчилна. Энд $R_{m,t}$ нь $t$-р хугацаан дахь
зах зээлийн өгөөж, $a_i$ нь регрессийн тогтмол
коэффициент, $\beta_i$ нь $i$-р хөрөнгийн өгөөжийн
зах зээлийн өөрчлөлтөд мэдрэг байдлыг илэрхийлэх
коэффициент, $\varepsilon_{ti}$ нь тухайн хөрөнгийн
үлдэгдэл өгөөж юм.

Нэг хүчин зүйлийн загварын регрессийн коэффициентийг
\[
\widehat{\beta}_i
=
\frac{
	\sum_{t=1}^{n}
	(y_{ti}-\overline{Y}_i)
	(R_{m,t}-\overline{R}_m)
}{
	\sum_{t=1}^{n}
	(R_{m,t}-\overline{R}_m)^2
}
\]
гэж үнэлнэ. Энд
\[
\overline{R}_m
=
\frac{1}{n}
\sum_{t=1}^{n}R_{m,t}
\]
нь зах зээлийн өгөөжийн түүврийн дундаж юм.

Регрессийн тогтмол коэффициентийг
\[
\widehat{a}_i
=
\overline{Y}_i
-
\widehat{\beta}_i\overline{R}_m
\]
гэж тооцно.

Тэгвэл $i$-р хөрөнгийн $t$-р хугацаан дахь
регрессийн үлдэгдэл
\[
\widehat{\varepsilon}_{ti}
=
y_{ti}
-
\widehat{a}_i
-
\widehat{\beta}_iR_{m,t}
\]
болно. Үүнийг
\[
\widehat{\varepsilon}_{ti}
=
(y_{ti}-\overline{Y}_i)
-
\widehat{\beta}_i
(R_{m,t}-\overline{R}_m)
\]
гэж илэрхийлж болно.

Нэг хүчин зүйлийн загварын үндсэн таамаглалаар
үлдэгдэл өгөөжүүд нь зах зээлийн хүчин зүйлтэй
хамааралгүй, мөн өөр өөр хөрөнгүүдийн үлдэгдэл
өгөөжүүд хоорондоо хамааралгүй байна. Өөрөөр хэлбэл,
\[
\operatorname{Cov}(R_m,\varepsilon_i)=0,
\qquad i=1,2,\ldots,N,
\]
ба
\[
\operatorname{Cov}(\varepsilon_i,\varepsilon_j)=0,
\qquad i\neq j.
\]

Эдгээр таамаглалын дор $i$ болон $j$-р хөрөнгийн
өгөөжийн ковариацыг
\[
\operatorname{Cov}(Y_i,Y_j)
=
\operatorname{Cov}
(a_i+\beta_iR_m+\varepsilon_i,
a_j+\beta_jR_m+\varepsilon_j)
\]
гэж илэрхийлнэ.

Ковариацын чанарыг ашиглавал
\[
\operatorname{Cov}(Y_i,Y_j)
=
\beta_i\beta_j\operatorname{Var}(R_m),
\qquad i\neq j,
\]
болно.

Харин $i$-р хөрөнгийн өгөөжийн вариац
\[
\operatorname{Var}(Y_i)
=
\beta_i^2\operatorname{Var}(R_m)
+
\operatorname{Var}(\varepsilon_i)
\]
болно.

Зах зээлийн өгөөжийн түүврийн вариацыг
\[
\widehat{\sigma}_m^2
=
\frac{1}{n}
\sum_{t=1}^{n}
(R_{m,t}-\overline{R}_m)^2
\]
гэж тооцно.

Мөн $i$-р хөрөнгийн регрессийн үлдэгдлийн
вариацыг
\[
\widehat{\sigma}_{\varepsilon,i}^{2}
=
\frac{1}{n}
\sum_{t=1}^{n}
\widehat{\varepsilon}_{ti}^{\,2}
\]
гэж тооцно.

Иймээс нэг хүчин зүйлийн зорилтот матриц
$T_{\mathrm{SF}}=(t_{ij})$-ын элементүүдийг
\[
t_{ij}
=
\begin{cases}
	\widehat{\beta}_i^2\widehat{\sigma}_m^2
	+
	\widehat{\sigma}_{\varepsilon,i}^{2},
	& i=j,\\[0.4em]
	\widehat{\beta}_i\widehat{\beta}_j
	\widehat{\sigma}_m^2,
	& i\neq j
\end{cases}
\]
гэж тодорхойлно \citep{ledoitwolf2003}.

Нэг хүчин зүйлийн зорилтот матрицыг матриц
хэлбэрээр
\[
T_{\mathrm{SF}}
=
\widehat{\boldsymbol{\beta}}
\widehat{\boldsymbol{\beta}}^{\top}
\widehat{\sigma}_m^2
+
\widehat D_{\varepsilon}
\]
гэж илэрхийлж болно.

Энд
\[
\widehat{\boldsymbol{\beta}}
=
(\widehat{\beta}_1,
\widehat{\beta}_2,
\ldots,
\widehat{\beta}_N)^{\top}
\]
нь регрессийн коэффициентүүдийн вектор, харин
\[
\widehat D_{\varepsilon}
=
\operatorname{diag}
\left(
\widehat{\sigma}_{\varepsilon,1}^{2},
\ldots,
\widehat{\sigma}_{\varepsilon,N}^{2}
\right)
\]
нь үлдэгдлийн вариацуудаас бүрдэх диагональ
матриц юм.

Нэг хүчин зүйлийн зорилтот матрицад суурилсан
шугаман агшилтын үнэлэгчийг
\[
\widehat{\Sigma}_{\mathrm{SF}}
=
(1-\widehat{\alpha}_{\mathrm{SF}})S
+
\widehat{\alpha}_{\mathrm{SF}}T_{\mathrm{SF}},
\qquad
0\leq\widehat{\alpha}_{\mathrm{SF}}\leq1,
\]
гэж тодорхойлно.

Энд $\widehat{\alpha}_{\mathrm{SF}}$ нь нэг хүчин
зүйлийн зорилтот матрицад харгалзах агшилтын
эрчим юм.

\paragraph{Агшилтын эрчмийн үнэлгээ.}

\citet{ledoitwolf2003} нь нэг хүчин зүйлийн
зорилтот матрицад суурилсан агшилтын эрчмийг
үнэлэхдээ түүврийн ковариацын матрицын
элементүүдийн түүвэрлэлтийн хэлбэлзэл,
түүврийн ковариацын матриц болон зорилтот
матрицын үнэлгээний алдаануудын хам хэлбэлзэл,
мөн эдгээр матрицуудын хоорондын зайг
харгалзан үзсэн.

Эхлээд түүврийн ковариацын матриц $S$ болон
нэг хүчин зүйлийн зорилтот матриц
$T_{\mathrm{SF}}$ хоорондын нормчилсон
квадрат Фробениусын зайг
\[
\widehat d_{\mathrm{SF}}
=
\left\|S-T_{\mathrm{SF}}\right\|^2=
\frac{1}{N}
\sum_{i=1}^{N}
\sum_{j=1}^{N}
(s_{ij}-t_{ij})^2
\]
гэж тодорхойлно.

Дараа нь түүврийн ковариацын матрицын
элементүүдийн түүвэрлэлтийн хэлбэлзлийг
нэгтгэн хэмжих хэмжигдэхүүнийг
\[
\widehat b_{\mathrm{SF}}
=
\frac{1}{nN}
\sum_{i=1}^{N}
\sum_{j=1}^{N}
\left[
\frac{1}{n}
\sum_{t=1}^{n}
\left\{
(y_{ti}-\overline{Y}_i)
(y_{tj}-\overline{Y}_j)
-s_{ij}
\right\}^{2}
\right]
\]
гэж тооцно.

Нэг хүчин зүйлийн зорилтот матриц нь
түүврийн өгөгдлөөс үнэлэгддэг учраас
$S$ болон $T_{\mathrm{SF}}$ матрицуудын
үнэлгээний алдаанууд хоорондоо
хамааралтай байж болно.

Иймээс агшилтын эрчмийг тооцоход
эдгээр үнэлгээний алдаануудын
хам хэлбэлзлийг харгалзах шаардлагатай.

Энэ хам хэлбэлзлийн нөлөөг
$\widehat c_{\mathrm{SF}}$ хэмжигдэхүүнээр
илэрхийлнэ.

Нэг хүчин зүйлийн зорилтот матрицын диагональ
элементүүд нь түүврийн ковариацын матрицын
диагональ элементүүдтэй тэнцүү буюу
$t_{ii}=s_{ii}$ байна.

\citet{ledoitwolf2003}-ийн Lemma 2-т
үндэслэн үнэлгээний алдаануудын
хам хэлбэлзлийг дараах байдлаар тооцно.

$i=j$ үед
\[
\widehat q_{ii}
=
\frac{1}{n}
\sum_{t=1}^{n}
\left\{
(y_{ti}-\overline{Y}_i)^2-s_{ii}
\right\}^{2}.
\]

Харин $i\neq j$ үед
\[
	\widehat q_{ij}
	=
	\frac{1}{n}
	\sum_{t=1}^{n}
	\Big[
	\widehat{\beta}_j
	(y_{ti}-\overline{Y}_i)
	(R_{m,t}-\overline{R}_m)
	+
	\widehat{\beta}_i
	(y_{tj}-\overline{Y}_j)
	(R_{m,t}-\overline{R}_m)
	-
	\widehat{\beta}_i\widehat{\beta}_j
	(R_{m,t}-\overline{R}_m)^2
	-t_{ij}
	\Big]\times
	\]
	\[\times
	\Big[
	(y_{ti}-\overline{Y}_i)
	(y_{tj}-\overline{Y}_j)-s_{ij}
	\Big].
\]

Иймээс хам хэлбэлзлийн нийлбэрийн
нормчилсон үнэлгээг
\[
\widehat c_{\mathrm{SF}}
=
\frac{1}{nN}
\sum_{i=1}^{N}
\sum_{j=1}^{N}
\widehat q_{ij}
\]
гэж тодорхойлно.
Ингэснээр нэг хүчин зүйлийн зорилтот
матрицад харгалзах агшилтын эрчмийг
\[
\widehat{\alpha}_{\mathrm{SF}}
=
\max
\left\{
0,
\min
\left\{
1,
\frac{
	\widehat b_{\mathrm{SF}}
	-
	\widehat c_{\mathrm{SF}}
}{
	\widehat d_{\mathrm{SF}}
}
\right\}
\right\}
\]
гэж үнэлнэ.

Энэ нь \citet{ledoitwolf2003}-ийн
асимптотик оновчтой агшилтын эрчмийн
үнэлгээг нормчилсон квадрат
Фробениусын зайд нийцүүлэн
илэрхийлсэн хэлбэр юм.

Энд $\widehat d_{\mathrm{SF}}>0$
байхыг шаардана.

Эцэст нь нэг хүчин зүйлийн зорилтот
матрицад суурилсан агшилтын ковариацын
матрицыг
\[
\widehat{\Sigma}_{\mathrm{SF}}
=
(1-\widehat{\alpha}_{\mathrm{SF}})S
+
\widehat{\alpha}_{\mathrm{SF}}T_{\mathrm{SF}}
\]
гэж олно.

\paragraph{Тооцоолох алгоритм.}

Нэг хүчин зүйлийн зорилтот матрицад
суурилсан агшилтын ковариацын
матрицыг дараах үе шаттайгаар тооцно.

\begin{enumerate}
	
	\item Зах зээлийн дундаж өгөөжийг тооцно.
	\[
	\overline{R}_m
	=
	\frac{1}{n}
	\sum_{t=1}^{n}R_{m,t}.
	\]
	

	
	\item Зах зээлийн өгөөжийн түүврийн
	вариацыг тооцно.
	\[
	\widehat{\sigma}_m^2
	=
	\frac{1}{n}
	\sum_{t=1}^{n}
	(R_{m,t}-\overline{R}_m)^2.
	\]
	
	\item Хөрөнгө тус бүрийн регрессийн
	коэффициентийг тооцно.
	\[
	\widehat{\beta}_i
	=
	\frac{
		\sum_{t=1}^{n}
		(y_{ti}-\overline{Y}_i)
		(R_{m,t}-\overline{R}_m)
	}{
		\sum_{t=1}^{n}
		(R_{m,t}-\overline{R}_m)^2
	}.
	\]
	
	\item Регрессийн тогтмол коэффициент
	болон үлдэгдлийг тооцно.
	\[
	\widehat a_i
	=
	\overline{Y}_i
	-
	\widehat{\beta}_i\overline{R}_m,
	\]
	\[
	\widehat{\varepsilon}_{ti}
	=
	y_{ti}
	-
	\widehat a_i
	-
	\widehat{\beta}_iR_{m,t}.
	\]
	
	\item Хөрөнгө тус бүрийн үлдэгдлийн
	вариацыг тооцно.
	\[
	\widehat{\sigma}_{\varepsilon,i}^{2}
	=
	\frac{1}{n}
	\sum_{t=1}^{n}
	\widehat{\varepsilon}_{ti}^{\,2}.
	\]
	
	\item Нэг хүчин зүйлийн зорилтот
	матрицыг байгуулна.
	\[
	t_{ij}
	=
	\begin{cases}
		\widehat{\beta}_i^2
		\widehat{\sigma}_m^2
		+
		\widehat{\sigma}_{\varepsilon,i}^{2},
		& i=j,\\[0.4em]
		\widehat{\beta}_i\widehat{\beta}_j
		\widehat{\sigma}_m^2,
		& i\neq j.
	\end{cases}
	\]
	
	Эдгээр элементүүдээр
	$T_{\mathrm{SF}}=(t_{ij})$
	матрицыг байгуулна.
	
	\item Түүврийн ковариацын матриц болон
	зорилтот матрицын хоорондын нормчилсон
	квадрат Фробениусын зайг тооцно.
	\[
	\widehat d_{\mathrm{SF}}
	=
	\frac{1}{N}
	\sum_{i=1}^{N}
	\sum_{j=1}^{N}
	(s_{ij}-t_{ij})^2.
	\]
	
	\item Түүврийн ковариацын матрицын
	элементүүдийн түүвэрлэлтийн
	хэлбэлзлийг тооцно.
	\[
	\widehat b_{\mathrm{SF}}
	=
	\frac{1}{nN}
	\sum_{i=1}^{N}
	\sum_{j=1}^{N}
	\left[
	\frac{1}{n}
	\sum_{t=1}^{n}
	\left\{
	(y_{ti}-\overline{Y}_i)
	(y_{tj}-\overline{Y}_j)
	-s_{ij}
	\right\}^{2}
	\right].
	\]
	
	\item Түүврийн ковариацын матриц
	болон нэг хүчин зүйлийн зорилтот
	матрицын үнэлгээний алдаануудын
	хам хэлбэлзлийг тооцно.
	
	$i=j$ үед
	\[
	\widehat q_{ii}
	=
	\frac{1}{n}
	\sum_{t=1}^{n}
	\left\{
	(y_{ti}-\overline{Y}_i)^2-s_{ii}
	\right\}^{2}.
	\]
	
	$i\neq j$ үед
\[
\widehat q_{ij}
=
\frac{1}{n}
\sum_{t=1}^{n}
\Big[
\widehat{\beta}_j
(y_{ti}-\overline{Y}_i)
(R_{m,t}-\overline{R}_m)
+
\widehat{\beta}_i
(y_{tj}-\overline{Y}_j)
(R_{m,t}-\overline{R}_m)
-
\widehat{\beta}_i\widehat{\beta}_j
(R_{m,t}-\overline{R}_m)^2
-t_{ij}
\Big]\times
\]
\[\times
\Big[
(y_{ti}-\overline{Y}_i)
(y_{tj}-\overline{Y}_j)-s_{ij}
\Big].
\]
	
	Дараа нь
	\[
	\widehat c_{\mathrm{SF}}
	=
	\frac{1}{nN}
	\sum_{i=1}^{N}
	\sum_{j=1}^{N}
	\widehat q_{ij}
	\]
	гэж тооцно.
	
	\item Агшилтын эрчмийг тооцно.
	
	$\widehat d_{\mathrm{SF}}>0$ үед
	\[
	\widehat{\alpha}_{\mathrm{SF}}
	=
	\max
	\left\{
	0,
	\min
	\left\{
	1,
	\frac{
		\widehat b_{\mathrm{SF}}
		-
		\widehat c_{\mathrm{SF}}
	}{
		\widehat d_{\mathrm{SF}}
	}
	\right\}
	\right\}.
	\]
	
	\item Агшилтын ковариацын
	матрицыг байгуулна.
	\[
	\widehat{\Sigma}_{\mathrm{SF}}
	=
	(1-\widehat{\alpha}_{\mathrm{SF}})S
	+
	\widehat{\alpha}_{\mathrm{SF}}T_{\mathrm{SF}}.
	\]
	
\end{enumerate}